\input{../header}
\chead{\textbf{Preview for 9/3}}

\begin{document}
%


%\onehalfspacing
\allowdisplaybreaks
%##################################################################
\def\arraystretch{3}
We computed a bunch of average velocities (that's what ``AV'' stands for, btw!) for the function $s(t) = 64-16(t-1)^2$ involving the point at $t=0.8$. Here are the ones where we went \textit{forward} from $t=0.8$:

{\large
\begin{tabular}{lllll}
    Time interval &=& $[0.8, 0.8+h]$ & AV setup & Result \\\hline
    $[0.8, 1.2]$ &=& $ [0.8, 0.8+(0.4)]$ &
    $\dfrac{s(1.2) - s(0.8)}{1.2 - 0.8}$ &
    0 ft/s \\
    $[0.8, 0.9]$ &=& $ [0.8, 0.8+\rule{1cm}{0.15mm}]$ &
    $\dfrac{s(\hspace{1cm}) - s(0.8)}{\rule{1cm}{0.15mm} - 0.8}$ &
    4.8 ft/s \\
    $[0.8, 0.81]$ &=& $ [0.8, 0.8+\rule{1cm}{0.15mm}]$ &
    $\dfrac{s(\hspace{1cm}) - s(0.8)}{\rule{1cm}{0.15mm} - 0.8}$ &
    6.24 ft/s \\
    $[0.8, 0.801]$ &=& $ [0.8, 0.8+\rule{1cm}{0.15mm}]$ &
    $\dfrac{s(\hspace{1cm}) - s(0.8)}{\rule{1cm}{0.15mm} - 0.8}$ &
    6.384 ft/s \\
\end{tabular}
}

\vspace{2em}

We also computed some going \textit{backward} from $t=0.8$, but also I'm going to write them funny because it doesn't matter what order you write the two times $t=a$ and $t=b$:

{\large
\begin{tabular}{lllll}
    Time interval &=& $[0.8, 0.8+h]$ & AV setup & Result \\\hline
    $[0.8, 0.4]$ &=& $ [0.8, 0.8+(-0.4)]$ &
    $\dfrac{s(0.4) - s(0.8)}{0.4 - 0.8}$ &
    12.8 ft/s \\
    $[0.8, 0.7]$ &=& $ [0.8, 0.8+\rule{1cm}{0.15mm}]$ &
    $\dfrac{s(\hspace{1cm}) - s(0.8)}{\rule{1cm}{0.15mm} - 0.8}$ &
    8 ft/s \\
    $[0.8, 0.79]$ &=& $ [0.8, 0.8+\rule{1cm}{0.15mm}]$ &
    $\dfrac{s(\hspace{1cm}) - s(0.8)}{\rule{1cm}{0.15mm} - 0.8}$ &
    6.56 ft/s \\
    $[0.8, 0.799]$ &=& $ [0.8, 0.8+\rule{1cm}{0.15mm}]$ &
    $\dfrac{s(\hspace{1cm}) - s(0.8)}{\rule{1cm}{0.15mm} - 0.8}$ &
    6.416 ft/s \\
\end{tabular}
}

\pagebreak

Then we did some algebra. We noticed that the setup of all of these average velocities is the same:
{\large
\[ AV_{[0.8,\  0.8+h]} = \dfrac{s(0.8 + h) - s(0.8)}{(0.8+h) - 0.8} \]
}

To figure out what $s(0.8 + h)$ was, we replaced the $t$ in $s(t) = 64 - 16(t-1)^2$ by $(0.8 + h)$:

\[s(0.8 + h) = 64 - 16\Big[ (0.8 + h) - 1 \Big]^2\]
Let's simplify this as much as possible. 

\begin{itemize}
\item First, simplify inside the $\big[\qquad \big]$ by combining the two numbers: $0.8 - 1 = $
\item Next, let's deal with the ``squared.'' Remember that $(\text{stuff})^2$ means $(\text{stuff}) \cdot (\text{stuff})$, and then use FOIL or the box method to multiply this out.
\item PEMDAS says our next thing is to multiply by \textbf{negative} 16. (Do you see why?) Distribute!
\item Finally, add 64.
\end{itemize}
Did you get that the most simplified version of $s(0.8+h) = 63.36 + 6.4h - 16h^2$?

\hrulefill

Where were we going with this, anyway? Oh right, we were computing $AV_{[0.8,\  0.8+h]}$. Replace $s(0.8 +h)$ in that big equation above by this simpler version.
{\large
\[ AV_{[0.8,\  0.8+h]} = \dfrac{\phantom{63.36 + 6.4h - 16h^2} - s(0.8)}{(0.8+h) - 0.8} \]
}
What else can you do in that big equation to simplify further?
{\large
\[ AV_{[0.8,\  0.8+h]} = \phantom{\dfrac{\phantom{63.36 + 6.4h - 16h^2} - s(0.8)}{(0.8+h) - 0.8}} \]
}

Notice that I kept the left-hand side of that equation the same, because that's our label that tells us what we are doing. 

\hrulefill

If you do legal algebra all the way through, you should end up with the following result:
{\large
\[AV_{[0.8,\  0.8+h]} = 6.4 - 16h\]
}

What does this \textit{mean}? I guess $6.4-16h$ is a recipe for calculating something, but calculating what? What's $AV_{[0.8,\  0.8+h]}$? An average velocity, apparently, but between what two times?
\end{document}